% Options for packages loaded elsewhere
\PassOptionsToPackage{unicode}{hyperref}
\PassOptionsToPackage{hyphens}{url}
\PassOptionsToPackage{dvipsnames,svgnames,x11names}{xcolor}
\documentclass[
]{article}
\usepackage{xcolor}
\usepackage[margin = 0.5in]{geometry}
\usepackage{amsmath,amssymb}
\setcounter{secnumdepth}{-\maxdimen} % remove section numbering
\usepackage{iftex}
\ifPDFTeX
  \usepackage[T1]{fontenc}
  \usepackage[utf8]{inputenc}
  \usepackage{textcomp} % provide euro and other symbols
\else % if luatex or xetex
  \usepackage{unicode-math} % this also loads fontspec
  \defaultfontfeatures{Scale=MatchLowercase}
  \defaultfontfeatures[\rmfamily]{Ligatures=TeX,Scale=1}
\fi
\usepackage{lmodern}
\ifPDFTeX\else
  % xetex/luatex font selection
\fi
% Use upquote if available, for straight quotes in verbatim environments
\IfFileExists{upquote.sty}{\usepackage{upquote}}{}
\IfFileExists{microtype.sty}{% use microtype if available
  \usepackage[]{microtype}
  \UseMicrotypeSet[protrusion]{basicmath} % disable protrusion for tt fonts
}{}
\makeatletter
\@ifundefined{KOMAClassName}{% if non-KOMA class
  \IfFileExists{parskip.sty}{%
    \usepackage{parskip}
  }{% else
    \setlength{\parindent}{0pt}
    \setlength{\parskip}{6pt plus 2pt minus 1pt}}
}{% if KOMA class
  \KOMAoptions{parskip=half}}
\makeatother
\usepackage{graphicx}
\makeatletter
\newsavebox\pandoc@box
\newcommand*\pandocbounded[1]{% scales image to fit in text height/width
  \sbox\pandoc@box{#1}%
  \Gscale@div\@tempa{\textheight}{\dimexpr\ht\pandoc@box+\dp\pandoc@box\relax}%
  \Gscale@div\@tempb{\linewidth}{\wd\pandoc@box}%
  \ifdim\@tempb\p@<\@tempa\p@\let\@tempa\@tempb\fi% select the smaller of both
  \ifdim\@tempa\p@<\p@\scalebox{\@tempa}{\usebox\pandoc@box}%
  \else\usebox{\pandoc@box}%
  \fi%
}
% Set default figure placement to htbp
\def\fps@figure{htbp}
\makeatother
\setlength{\emergencystretch}{3em} % prevent overfull lines
\providecommand{\tightlist}{%
  \setlength{\itemsep}{0pt}\setlength{\parskip}{0pt}}
\usepackage{color}
\usepackage{framed}
\setlength{\fboxsep}{.8em}

\usepackage{tcolorbox}

\tcbset{
  mystyle/.style={
    before skip=-10pt, % Adjust this value
    after skip=-10pt   % Adjust this value
  }
}
\newtcolorbox{mynoteblock}[1][]{mystyle, #1}

\newtcolorbox{blackbox}{
  colback=black,
  colframe=orange,
  coltext=white,
  boxsep=5pt,
  arc=4pt}
  
\newtcolorbox{ltbluebox}{
  colback=blue!5!white,
  colframe=blue!75!black,
  boxsep=5pt,
  arc=4pt}
  
\newtcolorbox{ltgreenbox}{
  colback=green!5!white,
  colframe=blue!75!black,
  boxsep=5pt,
  arc=4pt}
  
\newtcolorbox{ltpurplebox}{
  colback=purple!5!white,
  colframe=blue!75!black,
  boxsep=5pt,
  arc=4pt}
  


\newenvironment{cols}[1][]{}{}

\newenvironment{col}[1]{\begin{minipage}[t]{#1}\ignorespaces\vspace{0pt}}{%
\end{minipage}
\ifhmode\unskip\fi
\aftergroup\useignorespacesandallpars}

\def\useignorespacesandallpars#1\ignorespaces\fi{%
#1\fi\ignorespacesandallpars}

\makeatletter
\def\ignorespacesandallpars{%
  \@ifnextchar\par
    {\expandafter\ignorespacesandallpars\@gobble}%
    {}%
}
\makeatother

\usepackage{awesomebox}
\setlength{\aweboxvskip}{1mm}


\pagenumbering{gobble}

\usepackage{enumitem}
\setlist{itemsep=1pt, parsep=0pt} %# Adjust vertical spacing
\setlist[itemize,1]{leftmargin=*, labelsep=0.3em} %# Adjust horizontal spacing

\usepackage{bookmark}
\IfFileExists{xurl.sty}{\usepackage{xurl}}{} % add URL line breaks if available
\urlstyle{same}
\hypersetup{
  colorlinks=true,
  linkcolor={Maroon},
  filecolor={Maroon},
  citecolor={Blue},
  urlcolor={blue},
  pdfcreator={LaTeX via pandoc}}

\author{}
\date{\vspace{-2.5em}}

\begin{document}

\vspace{-1cm}

\section{MAFMC Indicators Project: Workshop 1
Overview}\label{mafmc-indicators-project-workshop-1-overview}

Workshop objective: Identify and prioritize Council management decisions
that could be supported or improved by operational use of ecosystem
indicators.

\subsection{Workshop Introduction}\label{workshop-introduction}

\begin{noteblock}

Presentations:

\begin{enumerate}
\def\labelenumi{\arabic{enumi}.}
\item
  A quick
  \href{https://www.mafmc.org/s/Project-5_IRA-Project-Brief.pdf}{project
  overview} for background
\item
  \emph{Definitions} of terms so we are all on the same page
\item
  Example indicator uses in management, followed by a preliminary list
  of example management decisions
\end{enumerate}

\end{noteblock}

\begin{ltgreenbox}

\subsection{Workshop Discussions and Desired
Outcomes}\label{workshop-discussions-and-desired-outcomes}

\textbf{Goal:} Which management decisions should this project focus on?
Why?

\textbf{Discussion 1:} Review the preliminary list of decisions

\textbf{Discussion 2:} Identify priority decisions

\textbf{Discussion 3:} Identify key information for priority decisions

\end{ltgreenbox}

\begin{cols}

\begin{col}{0.49\textwidth}

\emph{Things to think about\ldots{}}

\subsection{Discussion 1}\label{discussion-1}

\begin{itemize}
\tightlist
\item
  Are there Council decisions that should be added to the list of those
  presented?\\
\item
  Which decisions are unclear to you?\\
\item
  Can anything be removed?\\
\item
  \emph{Key Question: What management objectives are associated with
  each of the decisions?}
\end{itemize}

\subsection{Discussion 2}\label{discussion-2}

\begin{itemize}
\tightlist
\item
  Which management decisions most often require additional
  information?\\
\item
  What would a better decision look like? e.g., Reduce uncertainty?
  Increase consideration of tradeoffs? Prevent perverse incentives?
  Provide more flexibility or a wider range of options for
  participants?\\
\item
  Given the objectives for each decision type, would ecosystem
  information help make better decisions?\\
\item
  \emph{Key Question: Which management decisions should this project
  focus on? Why?}
\end{itemize}

\subsection{Discussion 3}\label{discussion-3}

\begin{itemize}
\tightlist
\item
  Which of the gathered indicators look useful to address the priority
  management decisions?\\
\item
  Are there priority decision points that don't appear to be supported
  by gathered indicators?\\
\end{itemize}

\end{col}

\begin{col}{0.02\textwidth}
~

\end{col}

\begin{col}{0.49\textwidth}

\emph{Project Steps: This workshop focuses on Management Objectives and
Decisions, but notes linkages to other steps}

\includegraphics[width=1\linewidth]{WK1_Overview_files/figure-latex/projectoutline-1}

\begin{ltbluebox}
\emph{This workshop addresses project goals:}\\
* Review and assess existing indicators and evaluate their utility to
support management objectives\\
* Identify and develop priority ecosystem and habitat indicators,
including associated targets and thresholds as appropriate, and identify
specific pathways for incorporation into Council management processes

\end{ltbluebox}

\end{col}

\end{cols}

\end{document}
